% concatenated from Small_2_Sets_and_Riesz_Sets_20260817.tex
\documentclass[12pt, twoside, leqno]{article}

\usepackage{amsmath,amsthm}
\usepackage{amssymb}
\bibliographystyle{plain}

%\pagestyle{headings}

\theoremstyle{definition}
\newtheorem{theorem}{Theorem}[section]
\newtheorem{corollary}[theorem]{Corollary}
\newtheorem{lemma}[theorem]{Lemma}
\newtheorem{proposition}[theorem]{Proposition}
\newtheorem{problem}[theorem]{Problem}
\newtheorem{definition}[theorem]{Definition}
\newtheorem{remark}[theorem]{Remark}
\newtheorem{remarks}[theorem]{Remarks}
\newtheorem{example}[theorem]{Example}
\newtheorem{examples}[theorem]{Examples}

\newtheorem{mainthm}[theorem]{Main Theorem}

\theoremstyle{definition}
\newtheorem{resl}{}
\newtheorem{points}{}[section]


\numberwithin{equation}{section}

%opening

	
\begin{document}
	\title{Small-2 Sets Are Riesz Sets }

\author{ A. To-Ming Lau and A. Ülger}
%\date{}
\maketitle
\thispagestyle{empty}

%This paper is dedicated to the memory of our friend and colleague, Professor Garth H. Dales.

\renewcommand{\thefootnote}{}
\footnote{}
\footnote{2010 \emph{Mathematics Subject Classification}:
	Primary 43A46, 43A45; Secondary 43A20, 46J10}
\footnote{\emph{Key words and phrases}: Arens regularity, small-2 sets, Riesz sets, Rosenthal sets, Fourier algebra, Fourier-Stieltjes algebra.} 
	\begin{abstract}
		Let $ G $ be a compact metrizable Abelian group, $ L^{1}(G) $ its group algebra and $ M(G) $ its measure algebra. For each proper subset $ E $ of the dual group $ \hat{G} $, let $ L^{1}_{E}(G)=\{f\in L^{1}(G):\hat{f}=0 \text{ on } \hat{G}\setminus E \}$ and $M_{E}=\{\mu\in M(G):\hat{\mu}=0 \text{ on }\hat{G}\setminus E\} $. If $ M_{E}(G)=L^{1}_{E}(G) $ then the set $ E $ is said to be a Riesz sets. If $ M_{E}(G)*M_{E}(G)\subseteq L_{E}^{1}(G) $ then $ E $ is said to be a small-2 set. The main results of this paper are the following:\vspace{2mm}
		
		1. Every small-2 set is a Riesz set.\vspace{2mm}
		
		2. The ideal  $ L^{1}_{E}(G) $ is Arens regular iff $ E $ is a Riesz set.\vspace{2mm}
		
	\noindent	Let $ A=L_{E}(G) $ and equip $ A^{**} $ with the first Arens product.\vspace{2mm} 
		
		(3). The centre of $ A^{**} $ is $ Z(A^{**})=A+N(A^{**}) $, where $ N(A^{**})=\{r\in A^{**}:rA^{**}=\{0\}\} $.\vspace{2mm}
		
	\noindent These results settle three long-standing open problems in this area. 
	\end{abstract}
\section{Introduction} Let $ G $ be a compact metrizable Abelian group and $ \Re $ one of the four spaces $ L^{1}(G) $, $ M(G) $, $ C(G) $ or $ L^{\infty}(G) $. For each proper subset $ E $ of the discrete group $ \hat{G} $, we set\vspace{2mm}

 $\Re_{E}(G)=\{f\in \Re: \hat{f}=0 \text{ on } \hat{G}\setminus E\} $.\vspace{2mm}  
 
 \noindent The space $ L^{1}_{E}(G) $ is a closed ideal of the group algebra $ L^{1}(G) $; and, conversely, every closed ideal of $ L^{1}(G) $ is of this form. The space $ M_{E}(G) $ is a $ \sigma (M(G),C(G)) $-closed ideal of the measure algebra $ M(G) $. The ideal $ M_{E}(G) $ is the dual space of the quotient space $ C(G)/C_{\hat{G}\setminus E} (G)$ of $ C(G) $. Our reference about the Banach algebras $ L^{1}(G) $ and $ M(G) $ is Rudin's classical book \cite{rudin_01}.\vspace{2mm}
  
In this paper our aim is to present the proofs of the three results stated in the abstract together with the proofs of some other results closely related to the problems studied in the paper.\vspace{2mm}

 In the paper \cite {Esmailvandi_Filali_Galindo}, the authors of that paper present a series of results closely related to the problems studied in the present paper. Although the cited paper contains quite a few good results, the results given in it are partial results and do not solve any of the problems we consider in this paper. In this paper we settle three long-standing open problems in this area and present complete answers to all the problems studied, except for the  problem asking whether there is an infinite discrete group $ G  $ for which the Fourier algebra $ A(G) $ of $ G $ is Arens regular.\vspace{2mm}
 
  Before explaining our approach, we give some information about the notions involved.\vspace{2mm}
 
\textbf{Riesz Sets}. Around 1916, the Riesz brothers, Frigyes Riesz and Marcel Riesz, while determining the dual space of the disk algebra $ A(D) $, proved the following important result.\vspace{2mm}

\textbf{Theorem}. Let $ \mathbb{T} $ be the unit circle group and $ \mu\in M(\mathbb{T}) $ a measure. Suppose that, for each negative integer $ n $, $ \hat{\mu}(n)=\int_{-\pi}^{\pi}e^{-in\theta}d\mu(\theta)=0 $. Then $ \mu $ is absolutely continuous with respect to the Lebesgue measure of $ \mathbb{T} $.\vspace{2mm}

\noindent This theorem, in the above notation, says that\vspace{2mm}

 $ M_{\mathbb{N}}(\mathbb{T})=L^{1}_{\mathbb{N}}(\mathbb{T}) $.\vspace{2mm}
 
 \noindent The notion of "Riesz sets" for an abstract compact Abelian group $ G $ has been introduced by Y. Meyer in the paper \cite  {Meyer_01}. Since then this notion has became a subject of its own. In \cite {lust_piq_01} and \cite {lust_piq_02} Lust-Piquard, among other results, proved that the Banach space $ L^{1}_{E}(G) $ has the Radon-Nikodym property iff $ E $ is a Riesz set. In \cite {Godefroy_01}  G. Godefroy has studied the Riesz sets in connection with Havin-Moorey theorem. There are many other works about the Riesz sets in the literature.\vspace{1mm}
 
   On the space $ 2^{\hat{G}} $ consider the product topology. In \cite {Tardivel_01} Tardivel proved that the set $ \Sigma $ of the Riesz subsets of $ \hat{G} $ is a non-Borel coanalytic set in  $ 2^{\hat{G}} $ so that practically it is impossible to determine all the Riesz subsets of $ \hat{G} $.\vspace{2mm}

\textbf{Small-2 Sets}. In 1938, N. Wiener and A. Wintner in \cite {Wiener Wintner}, and later on several other mathematicians, showed that there are singular measures $ \mu\in M(\mathbb{T}) $ such that the measure $ \mu*\mu $ is absolutely continuous with respect to the Lebesgue measure of $ \mathbb{T} $. As a kind of generalization of this result, in his 1965 paper \cite{glicksberg_01}  I. Glicksberg gave some conditions on a set $ E $ in the dual group of a locally compact Abelian group $ G $ implying that $ M_{E}(G)*M_{E}(G)\subseteq L^{1}_{E}(G) $. Later on such a set is called a "small-2 set". Several mathematicians (see e.g. the papers \cite {graham_01}  , \cite {pigno_01} , \cite{Yamaguchi} and the references there) studied the small-2 sets in connection with Riesz sets and asked whether every small-2 set is a Riesz set. As far as we know the subject, as of the day, no solution of this problem has appeared in the literature.\vspace{2mm}

\textbf{Arens Regular Ideals of $ L^{1}(G) $}. Let $ A $ be, say, a commutative Banach algebra. In 1952 R. Arens showed that the second dual of $ A^{**} $ can be made into a Banach algebra. In general the algebra $ A^{**} $ is not commutative; when it is commutative the algebra $ A $ is said to be Arens regular. At the very beginning of the subject in 1960's it is proved by Civin and Yood that the group algebra $ L^{1}(G) $ is Arens regular iff the group $ G $ is finite. In contrast to this result, the characterization of the Arens regular closed ideals of $ L^{1}(G) $ has turned out to be quite a difficult problem. As of the day, the only positive result known, in the case where $ G $ is compact and Abelian, is this: If the set $ E $ is a Riesz set then the ideal $ L_{E}^{1}(G) $, as a Banach algebra of its own, is Arens regular \cite {a_ulger_01}.\vspace{2mm} 

Now we give some explanation about our approach to the problems studied in this paper. Let $ G $ be a metrizable compact Abelian group. For a proper subset $ E $ of $ \hat{G} $, let $ A=L^{1}_{E}(G) $ and $ B=M_{E}(G) $.  We consider these ideals as Banach algebras of their own. To prove the main results of the paper we proceed as follows. We show that\vspace{2mm}

(1). The algebra $ B $ is Arens regular iff the algebra $ A $ is Arens regular iff $ A^{**}A^{**}\subseteq A $ iff $ B^{**}B^{**}\subseteq A $.\vspace{2mm}

(2). The algebra $ B $ is Arens regular iff, for each $ \mu\in B $, the operator $ L_{\mu}:B\rightarrow B $, defined by $ L_{\mu}(\lambda)=\mu*\lambda $, is weakly compact.\vspace{2mm}

(3). For $ \mu\in B $, the operator  $ L_{\mu}:B\rightarrow B $, defined by $ L_{\mu}(\lambda)=\mu*\lambda $, is weakly compact iff $ \mu\in A $.\vspace{2mm}

\noindent These results show that the algebra $ A $ is Arens regular iff $ A=B $. From this we deduce that the ideal $ L^{1}_{E}(G) $ is Arens regular iff $ E $ is a Riesz set.\vspace{2mm}

(4). For any set $ E\subseteq\hat{G} $, by \cite [Proposition 1]  {Zhang}, the algebra $ A=L^{1}_{E}(G) $ has a multiplier bounded approximate identity (=MBAI). Such a MBAI induces a bounded linear projection $ P:B^{**}\rightarrow B^{**} $ such that $P(B^{**})=P(A^{**})=P(B)$. We show that the inclusion $ B*B\subseteq A $ implies that, for $\mu\in B$, the element $P(\mu)$ is in the centre of $A^{**}$. From this we deduce that the algebra $B$ is an ideal in its second dual. Then using the fact that the algebra $L^{1}(G)$ has a bounded approximate identity we show that the algebra $A$ is Arens regular. This result combined with the above result show that every small-2 set is a Riesz set.\vspace{2mm}

(5) To determine the centre of the algebra $ A^{**} $, the main fact used is the fact that the algebra $ B $ is a dual algebra so that, each $ m $ in the centre of  $ A^{**} $ is of the form $ m=\mu +r $ in the decomposition $ B^{**}=B+X^{\perp} $, where $ X=C(G)/C_{\hat{G}\setminus E}(G) $. Using the known result that $ A^{**}m\subseteq A $, we conclude that $ \mu\in A $  and $ rA^{**}=\{0\} $. This determine $ Z(A^{**}) $.\vspace{2mm}

(6) We also characterize the Arens regular quotients of $ L^{1}(G) $. For a set $ E\subseteq\hat{G} $, the algebra $ L^{1}(G)/L^{1}_{\hat{G}\setminus E}(G) $ is Arens regular iff $ E $ is a Rosenthal set. That is, $ L^{\infty}_{E}(G)=C_{E}(G) $.\vspace{2mm}

(7) For an arbitrary discrete group $H$, as an application of the abstract results proved in Section 3, we prove that the Fourier algebra $A(H)$ is Arens regular iff $B_{r}(H)$, the reduced Fourier-Stieltjes algebra of $H$, is an ideal in its bidual iff $B_{r}(H)^{**}B_{r}(H)^{**}\subseteq A(H)$.\vspace{2mm}


In Section 2 we introduce our notation and terminology. In this section we also recall a few known results related to Arens regularity notion. In Section 3 we work with two abstract commutative semisimple Banach algebras $ A $ and $ B $ having the relevant properties of the ideals $ L_{E}^{1}(G) $ and $ M_{E}(G) $, respectively. For these abstract Banach algebras, we prove the analogues of the results stated above. Then in Section 4, taking $ A=L_{E}^{1}(G) $ and $ B=M_{E}(E) $, we obtain easily the main results of the paper. In Section 5 we consider Arens regularity of the Fourier algebra $ A(H) $ and the reduced Fourier-Stieltjes algebra $ B_{r}(H) $ of an arbitrary discrete group $ H $. The abstract results proved in Section 3 apply to these algebras as well and permit us to obtain the results noted in (7).\vspace{2mm}


\noindent The problem whether there exists an infinite discrete group $ H $ for which the algebra $ A(H) $ is Arens regular is still unsolved. Among the problems considered in this paper, this is the only problem that we were not able to settle. \vspace{2mm}

The paper is essentially self-contained. The main ingredients of the proofs are the standard results of the subject, Rosenthal's $ \ell^{1} $- Theorem and the existence of a MBAI in the ideals $ L^{1}_{E}(G) $. \vspace{2mm}

 \section{Notation And Preliminary Results}
In this section we are going to introduce the notation and the terminology that we use in the paper. We shall also recall a few known results that we need. Our notation and terminology are quite standard. Concerning Banach algebras our reference is G. Dales's book \cite {Dales_01}.\vspace{2mm}

For any Banach space $ X $, by $ X^{*}$ we denote its dual space. We always consider $ X $ as naturally embedded into its bidual $ X^{**} $. For $ x \in X $ and $ f \in X^{*} $, by $ < x,f >  $, and also by $ < f,x > $, we denote the natural duality between  $ X $ and $ X^{*} $. For any subspace $ Y $ of $ X $, by $ Y^{\perp} $ we denote the annihilator of $ Y $ in $X^{*}$. The weak-star topology of $ X^{*} $ is $ \sigma(X^{*},X) $; the weak topology of $ X^{*} $ is $ \sigma(X^{*},X^{**}) $. By $ X_{1} $ we denote the closed unit ball of $ X $. For a non-empty subset $ D $ of $ X $, by $ <D> $ we denote the closed subspace of $ X $ generated by $ D $.\vspace{2mm}

In the rest of this section $ A $ is a commutative semisimple Banach algebra.\vspace{2mm}

By $ \Phi_{A} $ we denote the Gelfand space of $ A $. For $ a\in A $, by $ \hat{a} $ the Gelfand transform of $ a $. If the ideal $ A_{c}=\{a\in A:\text{ the  support of  $\hat{a}$ is compact} \} $ is dense in $ A $ then the algebra $ A $ is said to be Tauberian. If the algebra $ A $ is regular then $ A $ is Tauberian iff every proper closed ideal of $ A $ is contained in $ Ker(\varphi) $ for some $ \varphi\in\Phi_{A}$. For a non-empty subset $ D $ of $ A $, $ DD $ is the set defined by $ DD=\{ab:a,b\in D\} $. In particular, $AA=\{ab:a,b\in A\}$. We recall here the following result.\vspace{2mm}

$\bullet$ If the algebra $ A $ is regular and Tauberian then $ \overline{AA}=A $ \cite [Proposition 3.1.10]{dales_ulger_01}.\vspace{2mm}

For $ a\in A $ and $ f\in A^{*} $, $ a\cdot f $ is the functional on $ A $ defined by the equality $ <a\cdot f,b>=<f,ab> $. Clearly $ ||a\cdot f||\leq ||f||.||a|| $.  For $ n\in A^{**} $, the functional $ n.f $ on $ A $ is defined by $ <n\cdot f,a>=<n,a.f> $. Clearly $ ||n\cdot f||\leq ||n||.||f|| $.\vspace{2mm} 

\textbf{Arens Product On $ A^{**} $}. The second dual $ A^{**} $ of $ A $ can be made into a Banach algebra as follows. For $ m,n\in A^{**} $, the product $ mn $ of $ m $ and $ n $ is defined by the equality\vspace{2mm}

$ <mn,f>=<m,n\cdot f> $.\vspace{2mm}

\noindent Equipped with this multiplication (which is called the first Arens multiplication), $ A^{**} $ is a Banach algebra and has $ A $ as a closed subalgebra.  The algebra $ A^{**} $ is in general not commutative; if it is commutative then the algebra $ A $ is said to be \textit{Arens regular}. Even if the algebra $ A^{**} $ is not commutative, since $ A $ is commutative, for $ a\in A $ and $ m\in A^{**} $, $ am=ma $.  Among many other places, in Chapter 2 of the book \cite {Dales_01} and in the memoir \cite {dales_lau_01}  the reader can find ample information about Arens regularity notion. \vspace{2mm} 

\textbf{ Weakly Almost Periodic Functionals On $ A $}. For $ f\in A^{*} $, consider the set $ O(f)=\{a\cdot f:a\in A_{1}\} $. If the set $ O(f) $ is relatively weakly compact in $ A^{*} $, then the functional $ f $ is said to be \textit{weakly almost periodic on $ A $}. We note that if $ f $ is weakly almost periodic on $ A $ then it is also weakly almost periodic on $ A^{**} $. That is, the set $ \{m\cdot f:m\in A^{**}_{1}\} $ is relatively weakly compact in $ A^{*} $, so in $ A^{***} $. The subspace $ WAP(A) $ of $ A^{*} $ consisting of the weakly almost periodic functionals on $ A $ is a closed subspace of $ A^{*} $. The equality $ WAP(A)=A^{*} $ occurs iff $ A $ is Arens regular.  We always have\vspace{2mm}

$ \overline{Span(\Phi_{A})}\subseteq <AA^{*}>\subseteq <A^{**}A^{*}> $.\vspace{2mm}  

\noindent$\bullet$ The algebra $ A $ is an ideal in its bidual $ A^{**} $ iff, for each $ a\in A $, the multiplication operator $ L_{a}:A\rightarrow A $, defined by $ L_{a}(b)=ab $, is weakly compact.\vspace{2mm}

\textbf{Dual Banach Algebras.} A commutative dual Banach algebra is a pair $ (B,X) $ consisting of a commutative Banach algebra $ B $ and a Banach space $ X $ such that $ B=X^{*} $ and such that the multiplication of $ B $ is $ \sigma(B,X) $ continuous in each variable when the other variable is kept fixed. This last condition is equivalent to saying that, for $ u\in B $ and $ f\in X $, the functional $ u\cdot f $ is in $ X $. For instance, for any compact Abelian group $ G $ and any subset $E$ of the dual group $\hat{G}$, the pair $ (M_{E}(G),C(G)/C_{\hat{G}\setminus E}(G)) $ is a commutative, semisimple dual Banach algebra.\vspace{2mm}

The other notions will be introduced as needed.
\section{General Results} In this section we presents a series of general results for an abstract commutative semisimple Banach algebra $ B $ and a closed ideal $ A $ of it. Both $ A^{**}$ and $ B^{**}$ are equipped with the first Arens multiplication, as defined above. In the next section we shall see that we can get easily the results stated in the abstract by taking $ B=M_{E}(G) $ and $ A=L^{1}_{E} (G)$. The results proved in this section also apply to the Fourier algebra $ A(H) $ and the reduced Fourier-Stieltjes algebra $ B_{r}(H) $ of a discrete group $ H $, as well as to the closed ideals of $ A(H) $. \vspace{2mm}

\noindent  In this section\vspace{1mm}

$\bullet$ \textbf{the letter $ A $ denotes a commutative and semisimple Banach algebra and $G$ a metrizable compact Abelian group}.\vspace{2mm}

\noindent We shall not repeat these assumptions.  By\vspace{2mm}

$ Z(A^{**})=\{m\in A^{**}:mn=nm \text{ for all $ n$ in } A^{**}\} $\vspace{2mm}

\noindent we denote the centre of $ A^{**} $. \vspace{2mm}

\noindent We start with a simple lemma.

\begin{lemma}
	Suppose that the algebra $ A $ is an ideal in its bidual.  Then, for $m\in Z(A^{**})$ and $ f\in A^{*} $, the functional $ m\cdot f $ is in the space  $\overline{Span(\Phi_{A})} $.
\end{lemma}
\begin{proof}
	Suppose, for a contradiction, that for some $ m\in Z(A^{**}) $ and $ f\in A^{*} $, the functional $ m\cdot f $ is not in $\overline{Span(\Phi_{A})}$. Then there is an $ n\in A^{**} $ such that $ <n,m\cdot f>=<nm,f>\ne 0 $ but $ <n,\varphi>=0 $ for all $ \varphi $ in $ \overline{Span(\Phi_{A})} $. So, since $ A $ is an ideal in its second dual and since $ A $ is semisimple, for each $ a\in A $, $ an=0 $. Taking a net $ (a_{i}) $ in $ A $ that converges to $ m $ in the weak-star topology of $ A^{**} $, we get that $ mn=0 $. As $ m\in Z(A^{**}) $, $ nm=mn $ so that $ nm=0 $. But then $ <n,m\cdot f>=0 $. This contradiction proves that the functional $ m\cdot f $ is in the space $\overline{Span(\Phi_{A})} $.	
\end{proof}
The following result is the first key result of the paper. This result is proved in the paper \cite {a_ulger_02} under slightly stronger hypotheses. The proof given below is quite different and simpler than the one given in \cite {a_ulger_02}. In  \cite {a_ulger_02} the proof is based on Cantor's Diagonal Method; here the main ingredient of the proof is Rosenthal's $\ell^{1}$- Theorem, which says that every bounded sequence in a Banach space $X$ has a weakly Cauchy subsequence iff $X$ does not contain an isomorphic copy of $\ell^{1}$ \cite [Chapter 11]{Diestel_01}. The reader will note that, for any set $E$ in $\hat{G}$, the algebra $A=L^{1}_{E}(G)$ satisfies all the hypotheses of the next theorem. \vspace{2mm}

\begin{theorem} Suppose that the algebra $ A $ is weakly sequentially complete, an ideal in its bidual and its Gelfand space $ \Phi_{A} $ is discrete. Then an element $ m $ of $ A^{**} $ is in $ Z(A^{**}) $ iff $ mA^{**}\subseteq A $ and $ A^{**}m\subseteq A $. 	
\end{theorem}
\begin{proof}
	Suppose first that, for an $ m\in A^{**} $, we have $ mA^{**}\subseteq A $ and $ A^{**}m\subseteq A $. That is, for each $ n $ in $ A^{**} $, the products $ mn $ and $ nm $ are in $ A $. Since, for each $ \varphi\in \Phi_{A} $, $ <mn,\varphi>=<m,\varphi> <n,\varphi> $, we see that, for each  $ \varphi\in \Phi_{A} $, $ <mn,\varphi>=<nm,\varphi>  $. Hence, since $ A $ is semisimple and since both $ mn $ and $ nm $ are in $ A $, we conclude that $ mn=nm $. This shows that $ m $ is in the centre of the algebra $ A^{**} $.\vspace{2mm}
	
	To prove the converse, take an $ m\in Z(A^{**}) $. We want to prove that $ mA^{**}\subseteq A $ and $ A^{**}m\subseteq A $. To prove these inclusions, we first note that, since $ A $ is an ideal in its second dual, for $ a\in A $, the product $ ma $ is in $ A $. Let us see then that the multiplication operator $ L_{m}:A\rightarrow A $, defined by $ L_{m}(a)=am $, is weakly compact on $ A $. To prove this, let $ (a_{n}) $ be a sequence in the unit ball of $ A $. Since the space $ \Phi_{A} $ is discrete, the Banach space $ C_{0}(\Phi_{A}) $ does not contain an isomorphic copy of $ \ell^{1} $ \cite {semadeni_01}. Hence, by Rosenthal's $ \ell^{1} $- Theorem applied in the space $C_{0}(\Phi_{A})$, the sequence $ (\hat{a_{n}}) $ has a subsequence, denoted in the same way, that is weakly Cauchy in the space $ C_{0}(\Phi_{A}) $. So, for each $ \varphi\in \Phi_{A} $, the sequence $ (<a_{n},\varphi>) $ converges. Since the sequence $ (a_{n}) $ is bounded, for each $ f\in\overline{Span(\Phi_{A})} $, the sequence $ (<a_{n},f>) $ converges. So, since for each $ f\in A^{*} $, by the preceding lemma, the functional $ m\cdot f $ is in the space $ \overline{Span(\Phi_{A})} $, the sequence $ (<a_{n},m\cdot f>) $ converges. This shows that the sequence $ (ma_{n}) $, which is in $ A $, is weakly Cauchy. Hence, since $ A $ is weakly sequentially complete, the sequence $ (ma_{n}) $ converges weakly in $ A $. This proves that the operator $ L_{m} $ is weakly compact on $ A $. So $ L_{m}^{**}(A^{**})\subseteq A $. As $ L_{m}^{**}(A^{**})=A^{**}m $, $ A^{**}m\subseteq A $. Since $ m\in Z(A^{**}) $, $ mA^{**}\subseteq A $ too. This completes the proof.
 \end{proof}
 Under the hypotheses of this theorem, the preceding theorem shows that $$A \text{ is Arens regular iff } A^{**}A^{**}\subseteq A. $$ As a simple example to this theorem, we note that the sequence space $ A=\ell^{1} $, considered as a Banach algebra with coordinate-wise multiplication, satisfies all the hypotheses of this theorem. Since this algebra is Arens regular, either by this theorem or by direct verification, we have the inclusion $ A^{**}A^{**}\subseteq A $.\vspace{2mm}

\begin{remarks}
	(1). The reader may wonder whether the hypothesis that " $ A $ is an ideal in $ A^{**} $" implies that $ \Phi_{A} $ is discrete. This is not true. Indeed, there are infinite dimensional reflexive commutative semisimple Banach algebras whose Gelfand spaces are connected. For an example of such Banach algebras, see \cite [Example 3.3] {dales_ulger_01}.\vspace{2mm}
	
	(2). However, if the algebra $ A $ has the following "separating ball property"\vspace{2mm}
	
	For each pair of distinct functionals $ \varphi $ and $ \psi $ in $ \Phi_{A} $, there is an $ a\in A_{1} $ such that $ <\varphi,a>=1 $ and $ <\psi,a>=0 $, \vspace{2mm}
	
	\noindent then  the hypothesis that "$ A $ is an ideal in $ A^{**} $" implies that $ \Phi_{A} $ is discrete, see \cite [Corollary 3.2] {a_ulger_03}.\vspace{2mm}
	
	(3). If, for each $ a\in A $, the multiplication operator $ L_{a}:A\rightarrow A $, $ L_{a}(b)=ab $, is compact then, by the spectral properties of the compact operators, $ \Phi_{A} $ is discrete.\vspace{2mm}
	
	(4). The argument used in the proof of the preceding theorem can be used to obtain the following kind of results:\vspace{2mm}
	
	Let $ U $ be a quotient algebra of $ L^{1}(G) $. Suppose that $ U $ is Arens regular. Since $ U $ has a bounded approximate identity, the algebra $ U^{**} $ has a unit element so that, by the above lemma, $ \overline{Span(\Phi_{U})}=U^{**} $. Then, from the proof of the preceding theorem we see that whenever a quotient algebra $ U $ of $ L^{1}(G) $ is Arens regular, the algebra $ U $ does not contain an isomorphic copy of $ \ell^{1}$. The converse of this result seems to be an open problem, closely related to "Rosenthal sets", see Section 4 below.
	
\end{remarks}
The following theorem is another key result of this paper.\vspace{2mm}

\begin{theorem} Let $ D $ be a commutative semisimple Banach algebra. Suppose that the inclusion $ D^{**}D^{**}\subseteq D $ holds. Then the algebra $ D^{**} $ is Arens regular and an ideal in its bidual. 	
\end{theorem} 
\begin{proof}
	We first show that the algebra $ D $ is Arens regular. To see this, we note that, as at the beginning of the proof of the preceding theorem, the inclusion $ D^{**}D^{**}\subseteq D $ implies that, for $ m,n $ in $ D^{**} $, $ mn=nm $ on $ \Phi_{D} $. Since the products $ mn $ and $ nm $ are in $ D $ and since $ D $ is semisimple, we conclude that $mn=nm $ so that the algebra $ D^{**} $ is commutative, and the algebra $ D $ is Arens regular.\vspace{2mm}
	
	To see that the algebra $ D^{**} $ also is Arens regular, it is enough to show that each functional $ g\in D^{***} $ is weakly almost periodic on $ D^{**} $. To verify this, take a $ g\in D^{***} $. We first note that, since $ D^{***} $ is of the form \vspace{3mm}
	
	$ D^{***}=D^{*}\oplus D^{\perp} $, \vspace{3mm}  
	
	\noindent $ g $ is of the form $ g=f+\xi $ for some $ f\in D^{*} $ and $ \xi\in D^{\perp} $. Here $D^{\perp}$ is the annihilator of $D\subseteq D^{**}$ in $D^{***}$. Since, for each $ m,n $ in $ D^{**} $, by hypothesis, the product $ mn $ is in $ D $,  $ n\cdot \xi=0 $, so that \vspace{3mm}
	
	$ \{m\cdot g:m\in D^{**}_{1}\}=\{m\cdot f:m\in D^{**}_{1}\} $. \vspace{3mm} 
	
	\noindent Since the algebra $ D $ is Arens regular, the functional $ f $ is weakly almost periodic both on $ D $ and $ D^{**} $. So the set $ \{m\cdot g:m\in D^{**}_{1}\} $ is relatively weakly compact in $ D^{*} $, so in $ D^{***} $. This proves that the algebra $ D^{**} $ is Arens regular.\vspace{2mm}
	
	To prove that $ D^{**} $ is an ideal in its second dual, let $ m\in D^{**} $. We want to show that the multiplication operator $ L_{m}:D^{**}\rightarrow D^{**} $, defined by $ L_{m}(n)=nm $, is weakly compact. Since $ D^{**}D^{**}\subseteq D $,  $ D^{**}m\subseteq D $. So, the operator $ R_{m}:D\rightarrow D $, defined by $ R_{m}(d)=dm $, is weakly compact. As $ L_{m}=R_{m}^{**} $, the operator $ L_{m} $ is weakly compact. So, the algebra $ D^{**} $ is an ideal in its bidual. This implies that the algebra $ D $ is also an ideal in its bidual.
\end{proof}
In the rest of this section\vspace{2mm} 

$\bullet$ \textbf{The pair $ (B,X) $ will denote a commutative semisimple dual Banach algebra such that $ A $ a closed ideal of $ B $ and such that $ A_{1} $ is $ \sigma(B,X) $-dense in $ B_{1} $.}\vspace{2mm} 

The hypothesis that $ A_{1} $ is $ \sigma(B,X) $-dense in $ B_{1} $ implies that $ X $ embeds naturally and isometrically into $ A^{*} $. By "naturally" we mean that, for $ a\in A $ and $ f\in X $,\vspace{3mm} 

$ <a,f>_{B,X}=<f,a>_{A^{*},A} $\vspace{3mm}

\noindent so that the $ <B,X> $ duality agrees with the $ <A^{*},A> $ duality on $ A $. We note that, for any set $E\subseteq \hat{G}$, the spaces $B=M_{E}(G)$ and $A=L^{1}_{E}(G)$ satisfy the above conditions. \vspace{2mm}

The next result shows that the inclusion $ B^{**}B^{**}\subseteq A $ holds iff the inclusion $ A^{**}A^{**}\subseteq A $ holds. One direction of this equivalence is clear since $ A^{**}\subseteq B^{**} $; we prove only the reverse direction.\vspace{2mm}

\begin{theorem} Suppose that $ A^{**}A^{**}\subseteq A $. Then $ B^{**}B^{**}\subseteq A\subseteq B $ so that the algebra $ B^{**} $ is Arens regular and an ideal in its bidual.	
\end{theorem} 
\begin{proof}
	Before proving the inclusion $ B^{**}B^{**}\subseteq A $, let us first prove that the inclusion $ A^{**}A^{**}\subseteq A $ implies that the algebra $ B $ is Arens regular and an ideal in its bidual.\vspace{2mm}
	
	\noindent Let $ j:X\rightarrow A^{*} $ be the natural injection, $ j(f)=f $. As noted above, $ j $ is an isometry. Then $ j^{*}:A^{**}\rightarrow B $ is a surjective Banach algebra homomorphism that extends the natural injection $ A\hookrightarrow B $. That is, for $ a\in A $,  $ j^{*}(a)=a $. It follows that $ B $ is isomorphic to the quotient algebra $ A^{**}/Ker(j^{*}) $ of $ A^{**} $. So, since by the preceding theorem, the algebra $ A^{**} $ is Arens regular and an ideal in its bidual,  the algebra $ B $ also is Arens regular and an ideal in its bidual.\vspace{2mm}
	
	Next let us see that $ B^{**}B^{**}\subseteq A $. Since $ (B,X) $ is a dual algebra,\vspace{2mm}
	
	 $ B^{**}=B\oplus X^{\perp} $.\vspace{2mm}
	 
	 \noindent Let $ \tilde{u}=u+\lambda $ and $ \tilde{v}=v+\mu $ be two elements of $ B^{**} $ in this decomposition. Then, \vspace{3mm}
	
	$ \tilde{ u}\tilde{v}=uv+u\mu+v\lambda+\lambda \mu $.\vspace{3mm}
	
	\noindent Since $ (B,X) $ is a dual algebra,  for $ f\in X $, the functionals  $ u\cdot f$  and $v\cdot f$ are in $ X $. Since $ B $ is an ideal in $ B^{**} $, the products $ u\lambda $ and $ v\mu $ are in $ B $. Hence, for each $ f\in X $, $ <u\lambda,f>=<\lambda,u\cdot f>=0 $. So $ u\lambda=0$. For the same reasons, $v\mu=0 $ too. To see that $ \lambda\mu=0 $ too, let $ (u_{i}) $ be a net in $ B $ that converges to $ \lambda $ in the weak-star topology of $ B^{**} $. Since, as just seen, $ u_{i}\mu=0 $, we conclude that $ \lambda\mu=0 $ too. So\vspace{2mm}
	
	$ \tilde{u}\tilde{v}=uv$.\vspace{2mm}
	
	\noindent Hence $ B^{**}B^{**}\subseteq BB $.\vspace{ 2mm}
	
	To see that $ BB\subseteq A $, let $ u,v $ be two elements in $ B $ with $ ||u||\leq 1 $. Let $ (a_{i}) $ be a net in $ A_{1} $ that converges in the weak-star topology of $ B $ to $ u $. Then, since $ B $ is an ideal in its bidual, the multiplication operator $ L_{v} $ is weakly compact on $ B $, so on $ A $. Then the net $ (a_{i}v) $, which is contained in $ A $, is relatively weakly compact. So it converges weakly to $ uv $, so that $ uv $ is in $ A $. Thus,\vspace{2mm}
	
	$ B^{**}B^{**}\subseteq BB\subseteq A $.
	\end{proof}
The next result connects the preceding theorem and Theorem 3.2.

\begin{theorem}
	Suppose that the algebra $ A $ is an ideal in its second dual, weakly sequentially complete and $ \Phi_{A} $ is discrete. Then the algebra $ B $ is Arens regular iff for each $ u\in B $, the operator $ L_{u}:B\rightarrow B $ defined by $ L_{u}(v)=uv $, is weakly compact.
\end{theorem}
\begin{proof}
	Suppose first that the algebra $ B $ is Arens regular. Then the algebra $ A $ is Arens regular since $ A $ is a closed ideal of $ B $. So, by Theorem 3.2, $ A^{**}A^{**}\subseteq A $. Then, by Theorem 3.5, for each $ u\in B $, the operator $ L_{u} $ is weakly compact on $ B $.\vspace{2mm}
	
	Conversely, suppose that, for each $ u\in B $, the operator $ L_{u}:B\rightarrow B $ defined by $ L_{u}(v)=uv $, is weakly compact. In the space $ B^{**}=B\oplus X^{\perp} $ take two elements $ \tilde{u}=u+\lambda $ and $ \tilde{v}=v+\mu $. Then, \vspace{3mm}
	
	$ \tilde{ u}\tilde{v}=uv+u\mu+v\lambda+\lambda \mu $.\vspace{3mm}
	
	\noindent Since $ B $ is an ideal in its bidual, as seen in the proof of the preceding theorem, $u\mu=v\lambda=\lambda \mu=0   $ so that \vspace{2mm}
	
		$ \tilde{ u}\tilde{v}=uv $.\vspace{2mm}
		
		\noindent This shows that $ B^{**}B^{**}\subseteq B $. This, by Theorem 3.4, implies that $ B^{**} $, so $ B $, is Arens regular.
\end{proof}
In the next theorem we resume the results obtained so far.

\begin{theorem} Suppose that the algebra $ A $ is an ideal in its second dual, weakly sequentially complete and $ \Phi_{A} $ is discrete. Then the following assertions are equivalent.\vspace{2mm}
	
	(1) The algebra $ A $ is Arens regular.\vspace{2mm}
	
	(2) $ A^{**}A^{**}\subseteq A $.\vspace{2mm}
	
	(3) $ B^{**}B^{**}\subseteq A $.\vspace{2mm}
	
	(4) The algebra $ B $ is Arens regular.\vspace{2mm}
	
	(5) For each $ u\in B $, the operator $ L_{u}:B\rightarrow B $ defined by $ L_{u}(v)=uv $, is weakly compact.
	
\end{theorem}
Suppose that, under the hypotheses of the preceding theorem, the algebra $ A $ is Arens regular. Then, for each $ u\in B $, the operator $ L_{u} $ is weakly compact on $ B $. Hence, to prove that $ A=B $, it is enough to show that, for $ u\in B $, whenever the operator $ L_{u} $ is weakly compact, $ u $ is in $ A $. In the present abstract setting we do not know whether this is true or not. In the next section, using the fact that the algebra $ L^{1}(G) $ has a bounded approximate identity, we shall see that $ u\in A $ whenever $ L_{u} $ is weakly compact.\vspace{2mm}
 
	\begin{remark}Suppose that the algebra $ A $ is regular (this is the case if $ \Phi_{A} $ is discrete) and Tauberian. Then, as noted in Section 2, $\overline{AA}=A $. So, in the case where $ BB\subseteq A $, since then $ AA\subseteq BB\subseteq A $, we see that  $ \overline{BB}=A $.
\end{remark}
Since $ A $ is an ideal in $ B $, the Gelfand space of the algebra $ B $ is of the form\vspace{3mm}

$ \Phi_{B}=\Phi_{A}\cup H $,\vspace{3mm} 

\noindent where $ H $ is the hull of $ A $. Now suppose that $ BB\subseteq A $. Then, since $ B $ is semisimple, this latter inclusion implies that $ H=\emptyset $ so that $ \Phi_{B}=\Phi_{A} $. If, in addition, the algebra $ B $ is Tauberian then $ B=A $. We record this as a proposition.

\begin{proposition} Suppose that $ BB\subseteq A $ and the algebra $ B $ is Tauberian. Then $ B=A $. 	
\end{proposition}
We continue with a couple more results. In the next result our aim is to understand when $ \Phi_{A}\subseteq X $ and $ \overline{Span(\Phi_{A})}=X $. We recall that $ X\subseteq A^{*} $ and $ X $ is closed in $ A^{*} $.

\begin{proposition}
	Suppose that $ A $ is an ideal in its bidual and $ \Phi_{A} $ separates the points of $ B $. Then $ \Phi_{A}\subseteq X $ and $ \overline{Span\Phi_{A}}=X $.
\end{proposition}
\begin{proof}
	Let us first see that $ \Phi_{A}\subseteq X $. To see this, take a $ \varphi\in \Phi_{A} $ and, for a contradiction, suppose that $ \varphi\notin X $. Then there is an $ m\in A^{**} $ such that $ <m,\varphi>=1 $ and $ m=0 $ on $ X $. Choose an $ a\in A $ such that $ <a,\varphi>=1 $. Then the product $ ma$ is in $A $, so in $ B $. Since $ (B,X) $ is a dual algebra, for each $ f\in X $, $a\cdot f$ is in $X$, and $ <ma,f>=<m,a\cdot f>=0 $. So $ ma=0 $. Since $ <ma,\varphi>=<m,\varphi>\ne 0 $, we have a contradiction. Hence $ \Phi_{A}\subseteq X $ and $ \overline{Span(\Phi_{A})}\subseteq X $.\vspace{2mm}
	
	To prove that $ X\subseteq \overline{Span(\Phi_{A})} $, suppose that there is an $ f\in X $ that is not in $ \overline{Span(\Phi_{A})} $. Then, since $ B=X^{*} $, there is a $ u\in B $ such that $ <u,f>=1 $ and $ u=0 $ on $ \Phi_{A} $. Since $ \Phi_{A} $ separates the points of $ B $, $ u=0 $. This contradiction proves that $ \overline{Span(\Phi_{A})}=X $. 
\end{proof}
We note that if $ BB\subseteq A $ then $ \Phi_{A} $ separates the points of $ B $. If, in addition, $ A $ is an ideal in its bidual, then $\overline{Span(\Phi_{A})}=X $.\vspace{2mm}

 We shall need the following lemma in the next section. Before  we note that, for $ f\in A^{*} $ and $ u\in B $, the functional  $ u\cdot f $ on the algebra $ A $ is defined by $ <a,u\cdot f>=<au,f> $. Since $ A $ is an ideal in $ B $, this definition makes sense and $ u\cdot f $ is in $ A^{*} $.
 
\begin{lemma}
	Let $ u\in B $ be such that the operator $ L_{u}:A\rightarrow A $, defined by $ L_{u}(a)=ua $, is weakly compact. Then, for each $ f\in A^{*} $, the functional $ u\cdot f $ is in $ \overline{Span(\Phi_{A})} $.
\end{lemma}
\begin{proof}
	For a contradiction, suppose that, for some $ f\in A^{*} $, the functional $ u\cdot f $ is not in $ \overline{Span(\Phi_{A})} $. Then there is some $m\in A^{**} $ such that $ <m,u\cdot f>\ne 0 $ but $ m $ vanishes on $ \Phi_{A} $. Since $ L_{u} $ is weakly compact on $ A $, the product $ mu $ is in $ A $ and vanishes on $ \Phi_{A} $. So, since $ A $ is semisimple, $ mu=0 $. Then $ <m,u\cdot f>=0 $. This contradiction proves that $ u\cdot f $ is in the space $ \overline{Span(\Phi_{A})} $.
\end{proof}
Here we note that since the unit ball of $ A $ is weak-star dense in the unit ball of $ B $, for an $ u\in B $, the operator $ L_{u} $ is weakly compact on $ A $ iff it is weakly compact as an operator from $ B $ into itself.

\subsection{The Case Where The Algebra $ A $ Has A Multiplier Bounded Approximate Identity} If the algebra $ A $ has a multiplier bounded approximate identity (=MBAI) we can get more precise information about the products of the elements of $ A^{**} $. In this subsection our aim is to establish a result that will led us in the next section to the solution of the "small-2 set-Riesz set" problem. 

\begin{definition} A net $ (e_{\alpha})_{\alpha\in I} $ in $ A $ is said to be a MBAI (=multiplier bounded approximate identity) if there is a $ \beta>0 $ such that, for each $ a\in A $, $ ||ae_{\alpha}||\leq\beta||a|| $ and $ ||ae_{\alpha}-a||\rightarrow 0 $.	
\end{definition}
As follows from a result proved in \cite [Proposition 1]{Zhang} by Zhang, every closed ideal of the algebra $ L^{1}(G) $ has a MBAI. A short direct proof of this result is also given in \cite{Esmailvandi_Filali_Galindo}. For the algebra $ A=L^{1}_{E}(G) $, $ \beta=1 $. For that reason, to simplify the notation, we shall assume that $||ae_{\alpha}||\leq||a|| $ for all $ a $ in $ A $. \\

\begin{lemma}
	Suppose that the algebra $ A $ is an ideal in its bidual, has a MBAI $ (e_{\alpha}) $ and $ \Phi_{A} $ separates the points of $ B $ . Then, for each $ u\in B $, $ ||e_{\alpha}u||\leq ||u||
	$ and $ e_{\alpha}u\rightarrow u $ in the weak-star topology $ \sigma (B,X) $ of $ B $.
\end{lemma}
\begin{proof}
Let $ u $ be an element in the unit ball $ B_{1} $ of $ B $. Since the unit ball of $ A_{1} $ of $ A $ is $ \sigma(B,X) $-dense in $ B_{1} $, there is a net $ (a_{i}) $ in $ A_{1} $ such that $ ||a_{i}||\leq ||u|| $ for all $ i $ and $ a_{i}\rightarrow u $ in the weak-star topology of $ B $. So, for all $ \alpha $ and $ i $,\vspace{2mm}

$ ||e_{\alpha}a_{i}||\leq ||a_{i}||\leq ||u|| $.\vspace{2mm}

\noindent Hence, \vspace{2mm}

$ ||e_{\alpha}u||\leq\liminf_{i}||e_{\alpha}a_{i}||\leq ||u|| $. \vspace{3mm}

\noindent Since the net $ (e_{\alpha}u) $ is bounded, it has a $\sigma(B,X)$- cluster point, say $ v $ in $ B $. Since, by Proposition 3.10, $ \Phi_{A}\subseteq X $ and since, for each $ \varphi\in \Phi_{A} $, $ <e_{\alpha},\varphi>\rightarrow 1 $, we see that, for $ \varphi\in \Phi_{A} $, $ <v,\varphi=<u,\varphi> $. So, since $ \Phi_{A} $ separates the points of $ B $, $ v=u $. This proves that $ u $ is the only weak-star cluster point of the bounded net $ (e_{\alpha}u) $ so that $ e_{\alpha}u\rightarrow u $ in the weak-star topology $ \sigma(B,X) $ of $ B $.  

\end{proof}
Let now $ (e_{\alpha})_{\alpha\in I} $ be a fixed MBAI in the algebra $ A $. For each $ \alpha\in I $, consider the bounded linear operator $ L_{\alpha}:B\rightarrow B $, defined by $ L_{\alpha}(u)=e_{\alpha}u $. Let $ L_{\alpha}^{**} $ be the second adjoint of $ L_{\alpha} $, considered as an operator on $ B^{**} $. The net $ (L_{\alpha}^{**})_{\alpha\in I} $ is a bounded net in the operator algebra $ \mathbb{B}(B^{**})=(B^{**}\hat{\otimes}B^{*})^{*} $. So, this net has a subnet, denoted again as $ (L_{\alpha}^{**}) $, that converges in the weak-star topology $ \sigma(\mathbb{B}(B^{**}),B^{**}\hat{\otimes}B^{*}) $ of the algebra $ \mathbb{B}(B^{**}) $ to some operator $ P\in \mathbb{B}(B^{**}) $. That is, for each $ \tilde{u}\in B^{**} $ and $ g\in B^{*} $,\vspace{2mm}

$ <P(\tilde{u}),g>=\lim_{\alpha}<e_{\alpha}\tilde{u},g> $\vspace{2mm}

\noindent so that\vspace{2mm}

$ P(\tilde{u})=\sigma(B^{**},B^{*})-\lim_{\alpha}e_{\alpha}\tilde{u} $.\vspace{2mm}

Since the product $e_{\alpha}\tilde{u}$ is in $A$, so $\sigma(B^{**},B^{*})-\lim_{\beta}e_{\beta}(e_{\alpha}\tilde{u})=e_{\alpha}\tilde{u}$, we see that $P^{2}=P$ so that the operator $P$ is a bounded projection on $B^{**}$. This is fully proved in the paper\cite {lau_ulger_01}, and also in the book \cite [Theorem 2.3.84] {dales_ulger_02}. Since $P(B^{**})\subseteq P(A^{**})$, $P$ is also a projection on $A^{**}$ and $P(B^{**})=P(A^{**})$. Since, for $\tilde{u}=u+\xi$ in $B^{**}=B\oplus X^{\perp}$, $e_{\alpha}\tilde{u}=e_{\alpha}u$, we see that $$P(B^{**})=P(B).$$ Hence $$P(A^{**})=P(B).$$ Let $X_{A}=\overline{Span(\Phi_{A})}$. As on can see easily, and as proved in \cite [Proposition 4.3] {lau_ulger_01}, the kernel of $P:A^{**}\rightarrow A^{**}$ is $X_{A}^{\perp}$, the annihilator of $X_{A}$ in $A^{**}$. Hence $$ A^{**}=P(A^{**})\oplus X_{A}^{\perp}.$$ Since $P(m)=w^{*}-\lim e_{\alpha}m$, for $m,n$ in $A^{**}$, $$P(m)P(n)=P(mn).$$ Also, since, for $a\in A$ and $u\in B$, $aP(u)=au$, for $m\in A^{**}$, we have: $mP(u)=mu$. We shall use these observations below.



\begin{theorem} Suppose that the algebra $ A $ is weakly sequentially complete, an ideal in its second dual, has a MBAI $ (e_{\alpha})_{\alpha\in I} $, $\Phi_{A}$ is discrete and $ BB\subseteq A $. Then the algebra $B$, so the algebra $A$, is Arens regular.
	
\end{theorem}
\begin{proof} Our first aim is to prove that, for $u\in B$, $P(u)$ is in the centre of the algebra $A^{**}$. This will be achieved in several steps. Then, from this result we deduce that the algebra $B$ is Arens regular. \vspace{2mm}
	
(1) Let us first see that $A^{**}A^{**}\subseteq Z(A^{**})$. To see this,	let $m,n$ be two elements of $A^{**}$. We first note that, since $ A^{**}=P(A^{**})\oplus X_{A}^{\perp}$, for some $\lambda,\mu\in X_{A}^{\perp}$,  $m=P(m)+\lambda$ and $n=P(n)+\mu$. Since $A$ is an ideal in its bidual,  $AX_{A}^{\perp}=\{0\}$, so  $A^{**}X_{A}^{\perp}=\{0\}$. Hence $$mn=m(P(n)+\lambda)=mP(n)=(P(m)+\mu)P(n)=P(m)P(n)+\mu P(n).$$ Since $P(A^{**})=P(B)$, $P(m)=P(u)$ and $P(n)=P(v)$ for some $u,v$ in $B$. Since $BB\subseteq A$, the product $uv$ is in $A$ so that $P(uv)=uv$. Hence $$P(n)P(m)=P(v)P(u)=P(vu)=vu.$$ This show that the product $P(m)P(n)$ is in $A$, so in $Z(A^{**})$. \vspace{1mm}

Next we going to show that the product $\mu P(n)$ is also in $Z(A^{**})$. To see this, let $r$ be an element of $A^{**}$. Since $r\mu P(n)=0$, it is enough to show that $\mu P(n)r=0$ too. Using the fact that $r=P(r)+\nu=P(w)+\nu$, for some $w\in B$ and $\nu\in X_{A}^{\perp}$, we see that $$\mu P(n)r=\mu P(v)(P(w)+\nu)=\mu vw=vw\mu=0,$$ since the product $vw$ is in $A$. This completes the proof that $A^{**}A^{**}\subseteq Z(A^{**})$. \vspace{2mm}

(2) The inclusion $A^{**}A^{**}\subseteq Z(A^{**})$ implies that all the possible products of any three elements $n,m,s$ of $A^{**}$ are the same. i.e. $$smn=snm=msn=mns=nms=nsm.$$ \vspace{2mm}

(3) The fact that any three elements of $A^{**}$ freely commute implies that, for each $m\in A^{**}$ and each $f\in A^{*}$, the functional $m.f$ is weakly almost periodic on $A$. As $WAP(A)=X_{A}$, $m.f\in X_{A}$. This in turn implies that, for each $\lambda$ in $X_{A}^{\perp}$, $$<\lambda,m.f>=0.$$ Hence, for each $\lambda \in X_{A}^{\perp}$ and each $m\in A^{**}$, $\lambda m=0$. \vspace{2mm}

(4) Let now $u\in B$ be an element of $B$. Since $P(u)\in A^{**}$, for each $f\in A^{**}$, the functional $P(u).f$ is in $X_{A}$. So, for each $\lambda$ in $X_{A}^{\perp}$, $\lambda P(u)=0$. Since, each $n\in A^{**}$ is of the form  $n=P(n)+\lambda=P(v)+\lambda$, where $v\in B$, we see that $nP(u)=P(v)P(u)=uv$. As $P(u)n=P(u)P(v)=uv$, we see that $$A^{**}P(u)\subseteq A \text{ and } P(u)A^{**}\subseteq A.$$ From these inclusions, by Theorem 3.2, we conclude that $P(u)\in Z(A^{**})$.\vspace{2mm}

(5) To complete the proof, we have to show that the algebra $B$ is Arens regular. Since, for $m\in A^{**}$ and $u\in B$, $mP(u)=mu$, we see that $$A^{**}u\subseteq A. $$ This shows that the multiplier $L_{u}:A\rightarrow A$, $L_{u}(a)=au$, is weakly compact on $A$. Since $A_{1}$ is $\sigma(B,X)$-dense in $B_{1}$, and since $L_{u}:B\rightarrow B$, $L_{u}(v)=vu$, is $\sigma(B,X)$- continuous on $B$, $L_{u}$ is also weakly compact on $B$ so that the algebra $B$ is an ideal in its bidual. Hence, by Theorem 3.7, the algebra $B$, so the algebra $A$, is Arens regular.


\end{proof}  
\begin{remark}  To underline the importance of the above theorem, we note here the following instance of the theorem. \vspace{2mm}
	
	Let, for some subset $E$ of the dual group $\hat{G}$, $A=L^{1}_{E}(G)$ and $B=M_{E}(G)$. The preceding theorem proves that, for any small-2 set $E$, the algebra $A=L^{1}_{E}(G)$ is Arens regular.
	
\end{remark}

 We shall return to the above theorem in the next section. 

\section{The Proofs Of The Main Results}
In this section we shall present the proofs of the results stated in the abstract and also the proofs of some closely related results. In this section we suppose that $ G $ is a compact metrizable Abelian group. Actually, using the technique of \cite[Lemma 2.3]{npu}, one can prove that  all the results presented in this section are valid without the hypothesis that the group $ G $ is metrizable but the proofs become more complicates. For the clarity of the proofs we have preferred to work with a metrizable compact group.\vspace{2mm}

What we have in this section that we did not have in the previous section is that the algebra $ L^{1}(G) $ has a bounded approximate identity (=BAI). We shall exploit this.\vspace{2mm}

Before starting with the proofs, we first lay out the background of the subject.\vspace{2mm}

 We define $ (L^{1}(G),L^{\infty}(G)) $- duality, for $ f\in L^{1}(G) $ and $ \varphi\in L^{\infty} $, as \vspace{3mm}
 
 $ <f,\varphi>=\int_{G}f(-t)\varphi(t)dt $,\vspace{3mm}
 
 \noindent where $ dt $ is the Haar measure of $ G $. With this duality, for $ f,g $ in $ L^{1}(G) $ and $ \varphi\in L^{\infty}(G) $,\vspace{3mm}
 
 $ <g,f*\varphi>=<g*f,\varphi> $.\vspace{3mm}

\noindent Since the group $ G $ is metrizable, the algebra $ L^{1}(G) $ has a sequential BAI $ (e_{n}) $ bounded by one. Let $ e $ be a weak-star cluster point in $ L^{1}(G) ^{**}$ of this sequence. This $ e $ is a right identity in $ L^{1}(G) ^{**}$. Since the mapping $ P:L^{1}(G)^{**}\rightarrow L^{1}(G)^{**} $, defined by $ P(m)=em $, is a projection, we have the decomposition\vspace{3mm}

$ L^{1}(G)^{**}=P(L^{1}(G)^{**})\oplus Ker(P) $.\vspace{3mm}

\noindent We denote this decomposition as\vspace{3mm}

$ L^{1}(G)^{**}=eL^{1}(G)^{**}\oplus(1-e)L^{1}(G)^{**} $.\vspace{3mm}

\noindent Using this $ e $, we can embed $ M(G) $ isometrically into $ eL^{1}(G)^{**} $ as follows:\vspace{2mm}

For $ \mu\in M(G) $, consider the multiplier $ L_{\mu}:L^{1}(G)\rightarrow L^{1}(G) $, defined by $ L_{\mu}(f)=f*\mu $. Then the second adjoint $ L_{\mu}^{**} $ of $ L_{\mu} $ applies $ L^{1}(G)^{**} $ into itself. The mapping \vspace{3mm}

$ j:M(G)\rightarrow  L^{1}(G)^{**} $\vspace{3mm}

\noindent defined by, $ j(\mu)=L_{\mu}^{**}(e) $, is an isometric Banach algebra isomorphism, $ j(M(G))=eL^{1}(G)^{**} $ and $ (1-e)L^{1}(G)^{**}=C(G)^{\perp} $. Here we note that $ eL_{\mu}^{**}(e)=L_{\mu}^{**}(e) $ since $ ee=e $.\vspace{2mm}

\noindent This same mapping $ j $, for any subset $ E $ of $ \hat{G} $, also embeds $ M_{E}(G) $ isometrically and isomorphically into $ L^{1}_{E}(G)^{**} $ and  $ j(M_{E}(G))=e L_{E}^{1}(G)^{**} $.  The projection $ P $ above induces the following decomposition of $ L^{1}_{E}(G)^{**} $. \vspace{3mm}
 
 $ L_{E}^{1}(G)^{**}=e L_{E}^{1}(G)^{**}\oplus (1-e) L_{E}^{1}(G)^{**} $.\vspace{3mm}
 
 \noindent With $A=L_{E}^{1}(G) $ and $X_{A}= \overline{Span(\Phi_{A})}$, we have\vspace{3mm}
 
  $ (1-e) L_{E}^{1}(G)^{**}=X_{A}^{\perp} $.\vspace{3mm}
  
  \noindent Here $ X_{A}^{\perp} $ is the annihilator of $ X_{A} $ in $ A^{**} $. Thus, up to an isometric isomorphism,\vspace{3mm}
 
 (4.a)\hspace{2mm}  $  L_{E}^{1}(G)^{**}=M_{E}(G)\oplus X_{A}^{\perp} $.\vspace{3mm}
 
 \textbf{4.1. The Arens Regular Ideals Of The Group Algebra $ L^{1}(G) $ }.\vspace{3mm}
 
 After these preliminaries, we can now prove the following critical lemma.
  
\begin{lemma}
	Let $ \mu\in M_{E}(G) $ and suppose that the multiplier operator $ L_{\mu}:L_{E}^{1}(G)\rightarrow L^{1}_{E}(G) $, defined by $ L_{\mu}(f)=f*\mu $, is weakly compact. Then $ \mu\in L_{E}^{1}(G) $.
\end{lemma}
\begin{proof} To simplify the notation, we put $ A=L^{1}_{E}(G) $ and $ B=M_{E}(G) $. It is clear that the algebra $ A $ satisfies all the hypotheses we used in the previous section. It is also clear that $ \Phi_{A} $ separates the points of $ B $ and $ A_{1} $ is weak-star dense in $ B_{1} $. Since $ A^{*}=L^{\infty}(G)/A^{\perp} $, by the Hahn-Banach Theorem, we can and do consider each $ \varphi $ in $ A^{*} $ as the restriction to $ A $ of a functional in $ L^{\infty}(G)=L^{1}(G)^{*} $.\vspace{1mm}
	
	 These being noted, by Lemma 3.13, for each $ \varphi\in A^{*} $, the functional $ \mu* \varphi $ is in $ \overline{Span(\Phi_{A})} $.  Let $e $ and $ w $ be two weak-star cluster points of the sequence $ (e_{n}) $ in $ L^{1}(G)^{**} $. We can and do consider $L_{\mu}$ as a mapping from $L^{1}(G)$ into itself, defined as $L_{\mu}(f)=f*\mu$. We want to show that $ L_{\mu}^{**}(e)=L_{\mu}^{**}(w) $. To this end, we first note that, since $ e=w $ on $ \Phi_{A} $, $ e=w $ on the space $ \overline{Span(\Phi_{A})} $.
	 Since, for any functional $ \varphi\in A^{*} $, the functional $ \mu*\varphi $ is in the space $ \overline{Span(\Phi_{A})} $,\vspace{3mm}
	 
	 $ <e,\mu*\varphi>=<w,\mu*\varphi> $.\vspace{3mm}
	 
	 \noindent Since $ L_{\mu}^{*}(\varphi)=\mu*\varphi $,\vspace{3mm}

	  $ <e,\mu*\varphi>=<L_{\mu}^{**}(e),\varphi>$ and $<w,\mu*\varphi>=<L_{\mu}^{**}(w),\varphi> $. \vspace{3mm}
	  
	\noindent From these equalities, which are valid for all $\varphi\in A^{*}$, since $L_{\mu}^{**}(e)$ and $L_{\mu}^{**}(w)$ are in $A^{**}$, we see that $ L_{\mu}^{**}(e)=L_{\mu}^{**}(w) $. This shows that the bounded sequence $ (e_{n}*\mu) $, which is in $ A $, has only one weak-star cluster point in $ A^{**} $. So, this sequence is weakly Cauchy in $ A $. Hence, since $ A $ is weakly sequentially complete, the sequence $ (e_{n}*\mu) $ converges weakly to some function in $ A $. On the other hand, since the sequence $ (e_{n}*\mu) $ converges to $ \mu $ in the weak-star topology of $ M(G) $,  the sequence $ (e_{n}*\mu) $ converges weakly to $ \mu $. Since the sequence $ (e_{n}*\mu) $ is in $ L^{1}_{E}(G) $ and $ L_{E}^{1}(G) $ is closed in $ M(G) $, we conclude that $ \mu\in L^{1}_{E}(G) $.    
\end{proof}
The next result is the first main result of the paper.
\begin{theorem} Let $ E $ be a subset of $ \hat{G} $. Then the ideal $ L_{E}^{1}(G) $ is Arens regular iff the set $ E $ is a Riesz set.
	
\end{theorem}
\begin{proof}
	Suppose that the ideal $ A=L^{1}_{E}(G) $ is Arens regular. Then, by Theorem 3.2, $ A^{**}A^{**}\subseteq A $. Hence, by Theorem 3.5, the algebra $ B=M_{E}(G) $ is Arens regular and is an ideal in its second dual so that, for each $ \mu\in M_{E}(G) $, the multiplication operator $ L_{\mu} $ is weakly compact on $ M_{E}(G) $, so on $ L^{1}_{E}(G) $. Then, by the preceding lemma, $ \mu\in A $ so that $ B=A $. Hence, the set $ E $ is a Riesz set. As said above, the converse of this result is known, proved in \cite {a_ulger_01}.
\end{proof}
\textbf{4.2. Small-2 Sets Are Riesz Sets}.\vspace{3mm}

 Let $ E $ be a small-2 set in $ \hat{G} $. Here we want to prove that the algebra $ M_{E}(G) $, equivalently, the algebra $L^{1}_{E}(G)$, is Arens regular so that $ E $ is a Riesz set. This is the second main result of the paper.\vspace{2mm}

\begin{theorem}
	Every small-2 set in $ \hat{G} $ is a Riesz set.
\end{theorem}
\begin{proof}
	Let $ E\subseteq \hat{G} $ be a small-2 set, $ A=L^{1}_{E}(G) $  and  $ B=M_{E}(G) $. Since $ E $ is a small-2 set, $ B*B\subseteq A $. As we have proved in Theorem 3.14, the inclusion $B*B\subseteq A$ implies that the algebra $A=L^{1}_{E}(G)$ is Arens regular. Hence, by the preceding theorem, $E$ is a Riesz set, and $M_{E}(G)=L^{1}_{E}(G)$.
	
\end{proof} 

\textbf{4.3. The Centre Of The Algebra $ L_{E}^{1}(G)^{**} $.}\vspace{3mm}

In this subsection, for an arbitrary subset $ E $ of $ \hat{G} $, we shall determine the centre of the Banach algebra $ L_{E}^{1}(G)^{**} $.\vspace{2mm}

Let $ \Re(\hat{G}) $ be the coset-ring of the group $ \hat{G}  $. This is the Boolean ring generated by the cosets of the subgroups of $ \hat{G} $. For a set $ E\subseteq\hat{G} $, the algebra $ A=L^{1}_{E}(G) $ has a BAI iff $ E\in \Re(\hat{G}) $ \cite {liu_roo_wang}. \vspace{2mm}

Now suppose that $ E $ is in the ring $ \Re(\hat{G}) $ so that the algebra $ A=L^{1}_{E}(G) $ has a BAI $ (d_{n}) $. Let $ m\in Z(A^{**}) $ and $ d_{1} $,  $ d_{2} $ be two $ \sigma(A^{**},A^{*}) $- cluster points in $ A^{**} $ of the sequence $ (d_{n}) $. Then, since $ m\in Z(A^{**}) $, $ d_{1}m=d_{2}m=m $. This shows that the only weak-star cluster point of the bounded sequence $ (d_{n}m) $ is $ m $. So the sequence $ (d_{n}m) $ is weakly Cauchy in $ A $ and converges weakly to $ m $. This proves that, when $ E\in \Re(\hat{G}) $, $ Z(A^{**})=A $.\vspace{2mm}

As seen above, for a Riesz set $ E $, the algebra $ A=L^{1}_{E}(G) $ is Arens regular so that  $ Z(A^{**})=A^{**} $. These are two extreme cases.\vspace{2mm}

Let now $ E $ be an arbitrary subset of $ \hat{G} $ and $ A=L^{1}_{E}(G) $. Below by $ N(A^{**}) $ we denote the following subspace\vspace{2mm}

$ N(A^{**})=\{r\in A^{**}:rA^{**}=\{0\}\} $\vspace{2mm}

\noindent of $ A^{**} $. We note that the equality $ rA^{**}=\{0\} $ implies that $ A^{**}r=\{0\} $ so that $ N(A^{**})\subseteq Z(A^{**}) $.\vspace{2mm}

 The next result is the third main result of the paper. \vspace{2mm}

\begin{theorem} Let $ E $ be a proper subset of $ \hat{G} $ and $ A=L^{1}_{E}(G) $. Then	$ Z(A^{**})=A+ N(A^{**}) $.
	
\end{theorem}
\begin{proof}
	The inclusion $(A+N(A^{**}))\subseteq Z(A^{**}) $ being clear, we prove the reverse inclusion. To this end, take an $ m\in Z(A^{**}) $. Put $ B=M_{E}(G) $ and $ X=C(G)/C_{\hat{G}\setminus E}(G) $. Then $ B^{**}=B+X^{\perp} $. Since $ A^{**}\subseteq B^{**} $, $ m $ is of the form\vspace{2mm}
	
	$ m=\mu +r$ \vspace{2mm}
	
	\noindent for some $ \mu \in B$ and $ r\in X^{\perp} $. Let $ (e_{n}) $ be a bounded approximate identity in $ L^{1}(G) $ and $ e $ a weak-star cluster point of it in $ L^{1}(G)^{**} $. Since $ e_{n}r $ is in $ A $ and $ r=0 $ on $ \Phi_{A} $, we conclude that $ e_{n}r=0 $ and $ e_{n}m=e_{n}*\mu $. From this it follows that \vspace{2mm}
	
	$ em=L_{\mu}^{**}(e) $.\vspace{2mm}
	
		\noindent Since $ m\in Z(A^{**}) $, by Theorem 3.2, $ A^{**}m\subseteq A $. As $ A^{**}em=Am $, $ A^{**}L^{**}_{\mu}(e)\subseteq A $. This shows that the multiplier $ L_{\mu}:A\rightarrow A $, defined by $ L_{\mu}(f)=f*\mu $, is weakly compact. Hence, by Lemma 4.1, $ \mu\in A $. So, with $ f=\mu $, $ m $ is of the form\vspace{2mm}
	
	$ m=f+r $. \vspace{2mm}
	
	\noindent Since both  $ m $ and $ f $ are in $ Z(A^{**}) $, $ r\in Z(A^{**}) $ too. As $ Ar=\{0\} $, $ A^{**}r=\{0\} $. Since $ r\in Z(A^{**}) $, $ rA^{**}=\{0\} $ too so that $ r\in N(A^{**}) $, and $ Z(A^{**})=A+N(A^{**}) $. 
	
	
\end{proof}

 \textbf{4.4. Arens Regular Quotients Of $ L^{1}(G) $}.\vspace{2mm}

 Above we have determined the Arens regular ideals of $ L^{1}(G) $. It is natural to wonder what are the Arens regular quotient algebras of $ L^{1}(G) $. In this subsection we shall determine the Arens regular quotient algebras of $ L^{1}(G) $.\vspace{2mm}

 In this section $ G $ is a compact Abelian group. For a subset $ E $ of $ \hat{G} $, we define the ideal $ k(E) $ as \vspace{2mm}
 
 $ k(E)=\{f\in L^{1}(G):\hat{f}=0 \text{ on } E\} $. \vspace{2mm}
 
 \noindent Note that $ k(E)=L^{1}_{\hat{G}\setminus E}(G) $. Note also that \vspace{2mm}
 
 $ (L^{1}(G)/k(E))^{*}=L^{\infty}_{E}(G) $.\vspace{2mm}
 
 \noindent We recall that\vspace{2mm}
 
 $\bullet$ The set $ E $ is said to be a Rosenthal set if $ L^{\infty}_{E}(G)=C_{E}(G) $.\vspace{2mm}
 
 The Rosenthal sets also have been studied by Lust Piquard and several other mathematicians. In \cite {lust_piq_01} and \cite {lust_piq_02} Lust-Piquard proves that $ E $ is a Rosenthal set iff the Banach space $ C_{E}(G) $ has the Radon-Nikodym property. Below we shall see that $ E $ is a Rosenthal set iff the quotient algebra $ L^{1}(G)/k(E) $ is Arens regular.\vspace{2mm}
 
  We recall the following known result \cite [Corollary 2.3.34]{dales_ulger_02}.
 
 \begin{lemma}
 	Let $ D $ be a Banach algebra and $ I $ a closed ideal of it. Then the quotient algebra $ D/I $ is Arens regular iff $ I^{\perp}\subseteq WAP(D) $.
 \end{lemma} 
We also recall that $ WAP(L^{1}(G))=C(G) $. Combining these two results we get the following theorem.\vspace{2mm}

\begin{theorem} The algebra $ L^{1}(G)/k(E) $ is Arens regular iff $ E $ is a Rosenthal set.	
\end{theorem} 
Thus,\vspace{2mm}
 
$\bullet$ The ideal $ L^{1}_{E}(G) $ is Arens regular iff $ E $ is a Riesz set; and\vspace{2mm} 

$\bullet$ The quotient algebra $ L^{1}(G)/k(E) $  is Arens regular iff $ E $ is a Rosenthal set.

\section{Arens Regularity Of The Fourier Algebra $ A(G) $ Of A Discrete Group $ G $.}
 For an arbitrary locally compact group $ G $, the Fourier algebra $ A(G) $, the Fourier-Stieltjes algebra $ B(G) $ and the reduced Fourier-Stieltjes algebra $ B_{r}(G) $ of $ G $ have been introduced by P. Eymard in his seminal paper \cite {Eymard_01}  in 1964. Since then these algebras have been intensively investigated. In the monograph \cite{Kaniuth_lau_01} the reader can find apple information about these algebras. Here we recall only the facts that the algebra $ A(G) $ is a commutative semisimple regular and Tauberian Banach algebra; its Gelfand space, via evaluation functionals, is $ G $. When $ G $ is commutative, via Fourier transform, the algebra $ A(G) $ is isometrically isometric to the group algebra $ L^{1}(\hat{G}) $ of the dual group $ \hat{G} $ of $ G $.\vspace{2mm}
 
 For an amenable locally compact group $ G $, in \cite{lau_wong_01}  Lau and Wong proved that the algebra $ A(G) $ is Arens regular iff $ G $ is finite. For an arbitrary locally compact group $ G $, Forrest \cite {b_forrest_01}  proved that if the algebra $ A(G) $ is Arens regular then $ G $ is discrete and its amenable subgroups are finite. For a discrete group $ G $ containing the free group $ F_{2} $, Losert proved in \cite {losert_01}, in contrast with the group algebra $ L^{1}(H) $ of a compact Abelian group $ H $, that the centre of the algebra $ A(G)^{**} $ is strictly larger than $ A(G) $. As of the day, as far as we know, it is not known whether there is an infinite discrete group $ G $ for which the algebra $ A(G) $ is Arens regular. In this section we present what we know about this problem.\vspace{2mm}
 
 Let $ G $ be the discrete group, $ A=A(G) $ and $ B=B_{r}(G) $. Then the pair $ (B_{r}(G),C^{*}_{r}(G)) $, where $ C_{r}^{*}(G) $ is the reduced group $ C^{*} $-algebra, is a commutative semisimple dual Banach algebra. The algebras $ A(G) $ and $ B_{r}(G) $ have all the properties we used in Section 3 about the algebras $ A $ and $ B $ there. A proof of the fact that the closed unit of $ A(G) $ is weak-star dense in the closed unit ball of $ B_{r}(G) $ can be found in \cite [Theorem 4.3.25]{dales_ulger_02}. The general results proved in Section 3 permit us to state the following result.
 
 \begin{theorem}
 	Let $ G $ be a discrete group, $ A=A(G) $ and $ B=B_{r}(G) $. Then the following assertions are equivalent.\vspace{2mm}
 	
 	(1) The algebra $ A $ is Arens regular.\vspace{2mm}
 	
 	(2) The algebra $ B $ is Arens regular.\vspace{2mm}
 	
 	(3) $ B^{**}B^{**}\subseteq A $.\vspace{2mm}
 	
 	(4) The algebra $ B $ is an ideal in its second dual.
 	
 \end{theorem} 
We do not know whether the inclusion $ BB\subseteq A $ implies that the algebra $ A $ is Arens regular. Similarity of this problem to the "small-2 set-Riesz set" problem is quite striking. This inclusion implies that $ \Phi_{B}=\Phi_{A} $. We do not know either whether there is an infinite discrete group $ G $ for which $ \Phi_{B}=\Phi_{A} $. If $ BB\subseteq A $ and $ B $ is Tauberian then, as noted above in Proposition 3.9, $ B=A $.\vspace{2mm}
  
  Without extra hypotheses on $G$, these are all we know about Arens regularity of the algebra $ A(G) $. The main difficulty in working with this algebra is that we do not have a bounded approximate identity around; and, although for certain groups $G$ such as  $SL(2,\mathbb{R})$, the algebra $A(G)$ has a MBAI, in general we do not have a MBAI in $ A(G) $ either. In the case where the algebra $A(G)$ has a MBAI, Theorem 3.14 permits us to state the following result.
  \begin{theorem} Suppose that the algebra $A(G)$ has a MBAI. Then the algebra $A(G)$ is Arens regular iff $B_{r}(G)B_{r}(G)\subseteq A(G)$.
  	
  \end{theorem}
  If the group $G$ is amenable then the algebra $A(G)$ has a BAI. The results obtained above about the closed ideals of the group algebra of a compact abelian group are also valid for the closed ideals of $A(G)$. We shall not repeat the proofs of them. Actually, the results obtained in Section 3 apply to the closed ideals of $A(G)$ without any changes.
  
  
  
  \vspace{4mm}
  
\textbf{This paper is submitted to JFA on August 7, 2024.} \vspace{3mm}


\begin{thebibliography}{999}


\bibitem {Dales_01} H.G. Dales, Banach Algebras and Automatic Continuity, Oxford Science publications, Clarendon press, Oxford, 2000.\vspace{2mm}

\bibitem {dales_lau_01} H.G. Dales and A. T.-M. Lau, The second duals of Beurling algebras Mem. Amer. Math. Soc. 177 (2005), no. 836, vi+191 pp.

\bibitem {dales_ulger_02}  H.G. Gales and A. Ülger, Banach function algebras, Arens regularity and BSE norms, Canadian Mathematical Society Books in Mathematics, Springer, 2024.\vspace{2mm}
 
 \bibitem {dales_ulger_01} H. G. Dales and A. Ülger, Approximate identities in Banach function algebras, Studia Math. 226, 155–187 (2015).\vspace{2mm}
 
\bibitem {Diestel_01} J, Diestel, Sequences and series in Banach spaces, Graduate Texts in Math. vol. 92, Springer-Verlag, New York 1984.\vspace{2mm}

\bibitem {Esmailvandi_Filali_Galindo} R. Esmailvandi, M. Filali and J. Galindo, Arens regularity of the ideals of  the group algebra of a compact group, Proc. Roy. Soc. Edinburgh Sect. A 155 (2025), no. 2, 718–734. \vspace{2mm}

\bibitem  {Eymard_01}  P. Eymard, L'algebre de Fourier d'un groupe locallement compact, Bull. Soc. Math. France 92 (1964),181-236.\vspace{2mm}

\bibitem {b_forrest_01}  B. Forrest, Arens regularity and discrete groups, Pacific J. Math. 151 no. 2 (1991), 217–227.\vspace{2mm}

\bibitem {glicksberg_01}  I. Glicksberg, Fourier-Stieltjes transform with small supports, Ill. Jour. Math. 12 (1965), 418-427.\vspace{2mm}

\bibitem {graham_01}  C. C. Graham, Fourier-Stieltjes transforms with small supports, Illinois J. Math., 18 (1974), 532-534.\vspace{2mm}

\bibitem {Godefroy_01}  G. Godefroy, Sous-espaces bien disposposes de L1. Applications, Trans. Amer. Math. Soc. 286 (1984), 227-249.\vspace{2mm}

\bibitem {npu}  N. Işık, J. Pym and A. Ülger, The second dual of the group algebra of a compact group, J. London Math. Soc. (2) 35 (1987), no. 1, 135–148.

\bibitem {Kaniuth_lau_01}  E. Kaniuth and A. T.-M. Lau, Fourier and Fourier–Stieltjes Algebras on Locally Compact groups, Mathematical Surveys and Monographs 231, American Mathematical Society, Providence, Rhode Island, 2018.\vspace{2mm}

\bibitem {lau_ulger_01}  A. T.-M. Lau and A. Ülger, A unified approach to the Arens regularity and related problems for a class of Banach algebras associated with locally compact groups. Quarterly J. Math. Oxford(2) 72, 1311–1327 (2021).\vspace{2mm}

\bibitem {lau_wong_01}  A. T.-M. Lau and J. C. S. Wong, Weakly almost periodic elements in $ L^{\infty}(G) $ of a locally compact group, Proc. Amer. Math. Soc. 107 (1989), 321–359.\vspace{2mm}

\bibitem  {liu_roo_wang}  T.S. Liu, A. van Rooij and J. Wang, Projections and approximate identities for ideals in group algebras, Trans. Amer. Math. Soc. 175 (1973) 469–482. \vspace{2mm}

\bibitem  {losert_01} V. Losert, The centre of the bidual of Fourier algebras (discrete groups), Bull. London Math. Soc. 48 no. 6 (2016), 968–976.\vspace{2mm}

\bibitem {lust_piq_01} F. Lust-Piquard, Ensembles de Rosenthal et ensembles de Riesz, C.R. Acad. Sc. Paris, t. 282, Series A (1976), 833-835.\vspace{2mm}

\bibitem {lust_piq_02} Lust-Piquard, Proprietes geometriques des sous-espaces invariantspar translation de L1(G) et C(G), Semin. Geom. Espaces Banach, Exposé No:26 (1977-78), Ecole Polytechnique.\vspace{2mm}

\bibitem  {Meyer_01} Y. Meyer, Spectres des mesures et mesures absolument continues, Studia Math. 30 (1968), 87-99.\vspace{2mm}

\bibitem {semadeni_01} A. Pełczy´nski and Z. Semadeni, Spaces of continuous functions. III. Spaces  $ C(\Omega) $ for $ \Omega $ without perfect subsets. Studia Math. 18, 211–222 (1959).\vspace{2mm}

\bibitem  {pigno_01} L. Pigno, Convolution products with small Fourier-Stieltjes transforms, Illinois J. Math., 19 (197577-78. \vspace{2mm}

\bibitem {Tardivel_01} V. Tardivel, Ensembles de Riesz, Trans. Amer. Math. Soc. 305 (1988), 167-174.\vspace{2mm}

\bibitem {rudin_01} W. Rudin, Fourier analysis on groups, John Wiley and Sons Inc., New York, (1990). \vspace{2mm}

\bibitem {a_ulger_01} A. Ülger, Characterizations of Riesz sets. Math. Scand. 108 (2011), 264–278.\vspace{2mm}

\bibitem {a_ulger_02} A. Ülger, Central elements of $ A^{**} $ for certain Banach algebras $ A $ without bounded approximate identities. Glasgow Math. J. 41 (1999), 369–377.\vspace{2mm}

\bibitem {a_ulger_03} A. Ülger, Some results about the spectrum of commutative Banach algebras under the weak topology and applications. Monatsh. Math. 121, 353–379 (1996).\vspace{2mm}

\bibitem {Wiener Wintner} N. Wiener and A. Wintner, Fourier-Stieltjes transforms and singular convolutions, Amer. J. Math. 60 (1938), 513-522.\vspace{2mm}

\bibitem {Yamaguchi} H. Yamaguchi, Remarks on Riesz sets, Hokkaido Math. J., 7 (1978), 328-335.\vspace{2mm}

\bibitem {Zhang} Y. Zhang,   Approximate identities for ideals of Segal algebras on a compact group. J. Funct.Anal. 191 (2002), 123–131. \vspace{5mm}



\end{thebibliography} \vspace{4mm}



\noindent
Antony To-Ming Lau (Retired professor)
Department of Mathematical and Statistical Sciences, University of Alberta, Edmonton, Canada T6G 2G1.\\

\noindent
E-mail: anthonyt@ualberta.ca\\

\noindent
Ali Ülger (Retired professor)
Department of Mathematics, Koç University, Sariyer, Istanbul 34450, Türkiye.\\

\noindent
E-mail: aulger@ku.edu.tr

 

	       
\end{document}
	
	
	
